% Bewertungspaket zum Gesamttest «Elektrizität» — Quelle des PDFs (python3 scripts/build-lp-pdf.py elektrizitaet).
% Ganze Punkte je Teilschritt; Werte und Folgefehler-Fälle mit python3 nachgerechnet (03.10.2026).
\documentclass[11pt]{article}
\usepackage{../lp-druck}
\renewcommand{\lpname}{Elektrizität}
\renewcommand{\lpteilgebiet}{6.2}
\renewcommand{\lpfuss}{Bewertungspaket}
\newcolumntype{P}{>{\centering\arraybackslash}p{10mm}}
\newenvironment{raster}{\par\smallskip\noindent\tabularx{\linewidth}{@{}>{\raggedright\arraybackslash}p{38mm}P X@{}}\toprule
  \small\textsc{Teilschritt} & \small\textsc{P} & \small\textsc{Musterlösung}\\\midrule}{\bottomrule\endtabularx\par}
\newcommand{\fehler}[1]{\par\smallskip{\small\textcolor{lprot}{\bfseries Typische Fehler:} #1}}
\newcommand{\E}{\,{\footnotesize\bfseries\color{lpgruen}(E)}}
\begin{document}
\lpkopf{Bewertungspaket zum Gesamttest}{}

\begin{lpachtung}{Erst lösen, dann öffnen}
Dieses Paket enthält die Musterlösung. Wer es vor dem Lösen liest, misst am Ende nur, wie gut das Abschreiben gelingt.
\end{lpachtung}

\section*{1 · So gehst du vor}
\begin{enumerate}[leftmargin=*,itemsep=2pt]
\item \textbf{Gesamttest lösen} — vollständig, mit Rechenweg, ohne dieses Paket.
\item \textbf{Fotografieren oder scannen}, gut lesbar, eine Seite pro Bild. \textbf{Kein Name} auf den Bildern.
  Handyfotos tragen oft unsichtbar Standort, Zeit und Gerät mit (Metadaten). Lade darum lieber einen Scan oder ein
  Bildschirmfoto des Fotos hoch, oder entferne vorher die Standortangaben.
\item \textbf{Bewerten lassen.} Ob du dafür eine KI nutzt und welche, entscheidest du selbst und in eigener Verantwortung —
  je nachdem, was dir aufgrund deines Alters und deiner persönlichen Situation erlaubt ist (Altersgrenzen und
  Nutzungsbedingungen des Anbieters, allenfalls Einverständnis deiner Eltern). Lade nur deine Lösung und dieses PDF hoch,
  keine persönlichen Daten, und füge den Auftrag aus Abschnitt~2 ein. Ohne KI kannst du dich mit Abschnitt~3 selbst bewerten.
\item \textbf{Rückmeldung prüfen.} Stimmt, was die KI aus deiner Handschrift gelesen hat? Eine KI kann sich irren —
  im Zweifel gilt das Raster in Abschnitt~3.
\item \textbf{Gezielt wiederholen} nach Abschnitt~4.
\end{enumerate}

\needspace{14\baselineskip}
\section*{2 · Auftrag an die KI}
\begin{tcolorbox}[colback=lplinie!25,colframe=lplinie,boxrule=0.5pt,arc=1mm,fontupper=\small]
Du bist Korrektorin oder Korrektor für Physik an einer Schweizer Berufsmaturitätsschule. Im Anhang findest du (1)~das Bewertungspaket zum Gesamttest «Elektrizität» mit Musterlösung und Punkteraster und (2)~Fotos meiner handschriftlichen Lösung.\par\medskip
Geh so vor:\par
1. Schreib zu jeder Aufgabe G1 bis G6 zuerst ab, was du in meiner Lösung liest — Rechenweg und Ergebnis. Was du nicht sicher lesen kannst, markierst du mit [unsicher]. Nicht raten.\par
2. Vergib die Punkte genau nach dem Raster im Bewertungspaket (Abschnitt 3), ganze Punkte je Zeile. Ziehe einen Fehler nur einmal ab: Rechne ich mit einem falschen Zwischenergebnis richtig weiter, gibt es die Folgepunkte. Die Zeilen «Typische Fehler» sagen, welche Rasterzeile bei diesem Fehler 0 Punkte gibt — sie sind kein zusätzlicher Abzug.\par
3. Andere richtige Lösungswege zählen voll. Gleichwertige Schreibweisen zählen voll: Dezimalpunkt oder Dezimalkomma, andere Vorsilbe ($0.133\;\text{A} = 133\;\text{mA}$), Zehnerpotenz oder ausgeschriebene Zahl, sinnvoll gerundete Werte (drei signifikante Stellen). Zeilen mit (E) sind Ergebnispunkte: Sie verlangen das Ergebnis \emph{mit richtiger Einheit}; ohne oder mit falscher Einheit gibt es diesen Punkt nicht. Alle anderen Punkte gibt es nur mit nachvollziehbarem Weg (Formel, dann Werte mit Einheiten).\par
4. Begründe jeden Abzug in einem Satz und nenne die Fehlerart (zum Beispiel «Minuten nicht in Stunden umgerechnet»). Schreib die Musterlösung nicht ab — zeig mir, wo mein Fehler liegt.\par
5. Gib zum Schluss eine Tabelle «Aufgabe | Punkte | Begründung», die Gesamtpunktzahl (von 25) und — nach der Selbsteinschätzung in Abschnitt 4 — welches Kapitel des Leitprogramms ich wiederholen soll. Keine Note.\par
6. Fehlt ein Foto oder ist eine Aufgabe unlesbar, sag es, statt sie zu bewerten.
\end{tcolorbox}

\needspace{16\baselineskip}
\section*{3 · Musterlösung und Punkteraster}
\textbf{Allgemein:} Ganze Punkte je Zeile. Ein Fehler wird nur einmal abgezogen (Folgefehler). Gleichwertige Wege und
Schreibweisen zählen voll. In der Spalte P: \textbf{(E)}~= Ergebnispunkt, nur mit richtiger Einheit; alle anderen Punkte nur mit
nachvollziehbarem Weg. «Typische Fehler» nennt, welche Zeile 0 Punkte gibt, und die Punkte, die bleiben.

\begin{lpaufgabe}{G1}{Ballon, $Q = -8.0\;\text{nC}$}{3 P}
\begin{raster}
Aufgenommen oder abgegeben & 1 & aufgenommen: negative Ladung heisst Elektronenüberschuss\\
Anzahl & 1\E & $n = \dfrac{|Q|}{e} = \dfrac{8.0 \cdot 10^{-9}\;\text{C}}{1.602 \cdot 10^{-19}\;\text{C}} \approx 5.0 \cdot 10^{10}$ Elektronen\\
Ladung der Haare & 1 & $+8.0\;\text{nC}$: Die Ladung wurde nur verschoben, die Summe bleibt null.\\
\end{raster}
\fehler{«abgegeben»: erste Zeile 0 $\to$ 2 von 3\,P.
$\text{nC}$ als $10^{-6}\;\text{C}$ gelesen ($5.0 \cdot 10^{13}$) oder $|Q| \cdot e$ gerechnet: Anzahl 0 $\to$ 2 von 3\,P.
Haare «neutral» oder «auch negativ»: dritte Zeile 0 $\to$ 2 von 3\,P.}
\end{lpaufgabe}

\begin{lpaufgabe}{G2}{E-Bike-Akku, $36\;\text{V}$, $14\;\text{Ah}$}{4 P}
\begin{raster}
Ladung mit Weg & 1 & $Q = I \cdot t = 14\;\text{A} \cdot 3600\;\text{s}$\\
Ladung & 1\E & $Q = 50\,400\;\text{C}$\\
Betriebsdauer & 1 & $t = \dfrac{Q}{I} = \dfrac{14\;\text{Ah}}{7\;\text{A}} = 2\;\text{h}$ (oder $\dfrac{50\,400\;\text{C}}{7\;\text{A}} = 7200\;\text{s}$)\\
Energie & 1\E & $W = U \cdot Q = 36\;\text{V} \cdot 14\;\text{Ah} = 504\;\text{Wh}$\\
\end{raster}
\fehler{$14 \cdot 60 = 840\;\text{C}$ (Stunden nur in Minuten): beide Ladungszeilen 0; Dauer und Energie richtig gerechnet zählen $\to$ 2 von 4\,P.
Energie als $36 \cdot 50\,400\;\text{J} = 1\,814\,400\;\text{J}$ ohne Umrechnung in Wh: richtig, aber nicht in der verlangten Einheit $\to$ Energiezeile 0, 3 von 4\,P.}
\end{lpaufgabe}

\begin{lpaufgabe}{G3}{Backofen, $230\;\text{V}$, $9.5\;\text{A}$}{5 P}
\begin{raster}
Leistung & 1\E & $P = U \cdot I = 230\;\text{V} \cdot 9.5\;\text{A} = 2185\;\text{W}$\\
Umrechnen & 1 & $2185\;\text{W} = 2.185\;\text{kW}$, $45\;\text{min} = 0.75\;\text{h}$\\
Energie & 1\E & $E = P \cdot t = 2.185\;\text{kW} \cdot 0.75\;\text{h} \approx 1.64\;\text{kWh}$\\
Kosten & 1\E & $\text{Kosten} = E \cdot \text{Preis} = 1.64\;\text{kWh} \cdot 0.27\;\text{CHF/kWh} \approx 0.44\;\text{CHF}$\\
Aussage korrigiert & 1 & Strom wird nicht verbraucht: Was hineinfliesst, fliesst zurück. Umgesetzt (und bezahlt) wird Energie.\\
\end{raster}
\fehler{$E = 2185 \cdot 45 = 98\,325$ «kWh» (nicht in kW und h umgerechnet): Umrechnen und Energie 0; die Kosten mit diesem Wert richtig gerechnet zählen (Folgefehler) $\to$ 3 von 5\,P.
Aussage nur mit «es heisst Stromstärke» korrigiert: letzte Zeile 0 $\to$ 4 von 5\,P.}
\end{lpaufgabe}

\begin{lpaufgabe}{G4}{Verlängerungskabel, Kupfer, $2 \times 20\;\text{m}$, $1.5\;\text{mm}^2$}{4 P}
\begin{raster}
Leiterlänge & 1 & hin und zurück: $l = 2 \cdot 20\;\text{m} = 40\;\text{m}$\\
Widerstand & 1\E & $R = \rho \cdot \dfrac{l}{A} = 0.017\;\dfrac{\Omega\,\text{mm}^2}{\text{m}} \cdot \dfrac{40\;\text{m}}{1.5\;\text{mm}^2} \approx 0.453\;\Omega$\\
Ansatz Spannung & 1 & $U = R \cdot I$\\
Spannungsabfall & 1\E & $U = 0.453\;\Omega \cdot 8\;\text{A} \approx 3.63\;\text{V}$\\
\end{raster}
\fehler{Nur $20\;\text{m}$ gerechnet ($R \approx 0.227\;\Omega$, $U \approx 1.81\;\text{V}$): Leiterlänge 0, Rest folgerichtig $\to$ 3 von 4\,P.
$R = \rho \cdot \frac{A}{l}$ (Bruch vertauscht): Widerstand 0, Spannung folgerichtig $\to$ 3 von 4\,P.}
\end{lpaufgabe}

\begin{lpaufgabe}{G5}{$R_1 = 60\;\Omega$, $R_2 = 120\;\Omega$, $U = 24\;\text{V}$}{6 P}
\begin{raster}
a) $R_\text{ges}$ in Reihe & 1\E & $R_\text{ges} = R_1 + R_2 = 180\;\Omega$\\
a) Stromstärke & 1\E & $I = \dfrac{U}{R_\text{ges}} = \dfrac{24\;\text{V}}{180\;\Omega} \approx 0.133\;\text{A} = 133\;\text{mA}$\\
a) Teilspannungen & 1 & $U_1 = I \cdot R_1 = 8\;\text{V}$, $U_2 = I \cdot R_2 = 16\;\text{V}$ (Probe: $8\;\text{V} + 16\;\text{V} = 24\;\text{V}$)\\
b) $R_\text{ges}$ parallel & 1\E & $R_\text{ges} = \dfrac{R_1 \cdot R_2}{R_1 + R_2} = \dfrac{60\;\Omega \cdot 120\;\Omega}{180\;\Omega} = 40\;\Omega$\\
b) Teilströme & 1 & $I_1 = \dfrac{U}{R_1} = \dfrac{24\;\text{V}}{60\;\Omega} = 0.4\;\text{A}$, $I_2 = \dfrac{24\;\text{V}}{120\;\Omega} = 0.2\;\text{A}$\\
b) Gesamtstrom & 1\E & $I = I_1 + I_2 = 0.6\;\text{A}$ (Probe: $\dfrac{24\;\text{V}}{40\;\Omega} = 0.6\;\text{A}$)\\
\end{raster}
\fehler{Parallel $\frac{1}{R_\text{ges}} = \frac{1}{60\;\Omega} + \frac{1}{120\;\Omega} = 0.025\;\frac{1}{\Omega}$ und dann $R_\text{ges} = 0.025\;\Omega$ (Kehrwert vergessen): b)~$R_\text{ges}$ 0; Teilströme und Gesamtstrom aus $I_1 + I_2$ zählen $\to$ 5 von 6\,P.
In Reihe die Spannung halbiert ($12\;\text{V}$ und $12\;\text{V}$): Teilspannungen 0 $\to$ 5 von 6\,P.}
\end{lpaufgabe}

\begin{lpaufgabe}{G6}{Defektes Gehäuse, $R_\text{K} = 1.2\;\text{k}\Omega$}{3 P}
\begin{raster}
Körperstrom & 1\E & $I_\text{K} = \dfrac{U}{R_\text{K}} = \dfrac{230\;\text{V}}{1200\;\Omega} \approx 0.192\;\text{A} = 192\;\text{mA}$\\
FI & 1 & Er vergleicht Hin- und Rückstrom; $192\;\text{mA}$ fehlen auf dem Rückweg, weit über $30\;\text{mA}$ (mehr als das Fünffache): Er muss auslösen, spätestens nach $40\;\text{ms}$. (Die Zeitangabe ist nicht verlangt.)\\
Leitungsschutzschalter & 1 & Er bleibt ein: $0.192\;\text{A}$ sind weit unter $13\;\text{A}$ — er schützt die Leitung, nicht den Menschen.\\
\end{raster}
\fehler{$I_\text{K} = 230\;\text{V} \cdot 1200\;\Omega$: Körperstrom 0; die Begründungen zählen, wenn sie mit dem eigenen Wert folgerichtig sind $\to$ höchstens 2 von 3\,P.
«Der LS löst aus, weil der Strom gefährlich ist»: LS-Zeile 0 $\to$ 2 von 3\,P.}
\end{lpaufgabe}

\needspace{14\baselineskip}
\section*{4 · Selbsteinschätzung}
\begin{tabularx}{\linewidth}{@{}l X@{}}\toprule
22 – 25 P & Die geprüften Teile sitzen. Wo du Punkte verloren hast: das Kapitel dieser Aufgabe nochmals (Zuordnung unten).\\
17 – 21 P & Den schwächsten Teil nochmals: Simulation und Übungen des Kapitels, in dem du die meisten Punkte verloren hast.\\
11 – 16 P & Zurück zu den Kapiteln aller Aufgaben, in denen du Punkte verloren hast.\\
0 – 10 P & Zurück zu Kapitel 1 und von dort der Reihe nach weiter.\\\midrule
\multicolumn{2}{@{}p{\linewidth}@{}}{\small\textbf{Aufgabe $\to$ Kapitel} (gilt für alle Bereiche): G1 $\to$ Kapitel 1 · G2 $\to$ Kapitel 1 (Ladung) und 2 (Energie) · G3 $\to$ Kapitel 2 · G4 $\to$ Kapitel 3 · G5 $\to$ Kapitel 4 · G6 $\to$ Kapitel 5}\\\bottomrule
\end{tabularx}
\end{document}
